\chapter{Creio na Santa Igreja Católica}
\chaptermark{Creio na Santa Igreja Católica}

A Igreja ocupa um lugar central no Credo. Depois de confessar a fé em Deus Pai, no Filho e no Espírito Santo, a Igreja professa a sua fé na Santa Igreja Católica. A Igreja não é uma mera instituição humana, mas um mistério de salvação que tem a sua origem no desígnio eterno da Santíssima Trindade. O Catecismo da Igreja Católica, nos parágrafos 748 a 975, apresenta com profundidade este artigo de fé, mostrando que a Igreja é inseparável de Cristo e do Espírito Santo.

\section{Introdução (748--750)}

``Cristo é a luz da humanidade; e é, por conseguinte, o desejo ardente deste sagrado Concílio, reunido no Espírito Santo, que, anunciando o Evangelho a toda a criatura, leve a todos os homens a luz de Cristo que resplandece visivelmente na Igreja''. Estas palavras da Constituição dogmática \emph{Lumen Gentium} mostram que a Igreja recebe toda a sua luz de Cristo. Ela é como a lua que reflecte a luz do sol. A fé na Igreja depende inteiramente da fé em Cristo e no Espírito Santo. A Igreja é o lugar onde o Espírito Santo floresce e comunica a santidade.

A Igreja é uma realidade visível e espiritual, humana e divina. A sua origem está no desígnio eterno do Pai, realizado pelo Filho e revelado pelo Espírito Santo. Professar ``Creio na Igreja'' não significa crer nela como se fosse Deus, mas reconhecer que Deus mesmo a fundou e a sustenta como instrumento da salvação.

\section{A Igreja no desígnio de Deus (751--780)}

A palavra ``Igreja'' (do grego \emph{ekklesia}) significa ``convocação''. Designa a assembleia do povo de Deus. No Antigo Testamento, refere-se à assembleia do povo de Israel no Sinai. Os cristãos reconheceram-se como herdeiros desta convocação: Deus convoca agora todos os homens, de todas as nações, para formar o novo Povo de Deus.

A Igreja foi prefigurada desde o princípio do mundo, preparada na Antiga Aliança, instituída por Cristo e manifestada pelo Espírito Santo. O Pai eterno decidiu reunir numa só Igreja todos os que acreditam em Cristo. Esta Igreja foi preparada de modo admirável na história do povo de Israel. Cristo fundou a Igreja pregando o Evangelho do Reino e escolhendo os Doze com Pedro à cabeça. A Igreja nasce do lado aberto de Cristo na Cruz, da qual brotam sangue e água, símbolos dos sacramentos.

O Espírito Santo, enviado no Pentecostes, manifesta publicamente a Igreja e a envia em missão a todas as nações. A Igreja é missionária por natureza. Ela é o sacramento universal da salvação, sinal e instrumento da união dos homens com Deus e entre si.

A Igreja é simultaneamente visível e espiritual, sociedade hierárquica e Corpo Místico de Cristo. Ela é ``uma realidade complexa que resulta de um elemento humano e de um elemento divino''. A sua estrutura visível está ao serviço da realidade espiritual.

\section{A Igreja --- povo de Deus, Corpo de Cristo, templo do Espírito Santo (781--810)}

A Igreja é o Povo de Deus. Deus não pertence a nenhum povo em particular, mas adquiriu para si um povo de entre aqueles que antes não eram povo. Torna-se membro deste povo não por nascimento carnal, mas pelo novo nascimento da água e do Espírito, pela fé em Cristo e pelo Baptismo. Este Povo tem como cabeça Cristo, a sua lei é o mandamento novo do amor e a sua missão é ser sal da terra e luz do mundo.

O Povo de Deus participa dos três ofícios de Cristo: sacerdotal, profético e real. Todo o baptizado participa, segundo a sua condição, da função sacerdotal, profética e régia de Cristo. Jesus é o Ungido do Espírito; o seu Povo participa desta unção.

A Igreja é também o Corpo de Cristo. Cristo é a Cabeça e os baptizados são os membros. Esta imagem manifesta a unidade orgânica da Igreja e a diversidade de funções. O Espírito Santo é a alma deste Corpo, que lhe dá vida e unidade.

A Igreja é igualmente o Templo do Espírito Santo. O Espírito habita na Igreja e nos corações dos fiéis como num templo. Ele edifica a Igreja através dos sacramentos, da Palavra e dos carismas. A Igreja é o lugar onde o Espírito Santo distribui os seus dons para a utilidade comum.

\section{A Igreja é Una, Santa, Católica e Apostólica (811--870)}

``Esta é a única Igreja de Cristo, que no Credo professamos ser una, santa, católica e apostólica''. Estas quatro características, inseparavelmente ligadas, indicam traços essenciais da Igreja. Elas não são obra da Igreja, mas de Cristo que, pelo Espírito Santo, a torna una, santa, católica e apostólica.

A Igreja é \textbf{una} porque tem uma única origem (a Trindade), um único Fundador (Cristo) e uma única alma (o Espírito Santo). A unidade é da sua essência. A Igreja é \textbf{santa} porque Cristo a santificou e lhe deu o Espírito Santo como fonte de santidade. A Igreja é \textbf{católica} (universal) porque está enviada a toda a humanidade e contém em si a plenitude dos meios de salvação. A Igreja é \textbf{apostólica} porque foi fundada sobre os Apóstolos, conserva a sua doutrina e continua a ser guiada pelos sucessores dos Apóstolos.

Estas marcas são visíveis, embora realizem-se de forma imperfeita na história. A Igreja peregrina caminha para a perfeição que só receberá na glória do céu.

\section{Os fiéis de Cristo --- Hierarquia, leigos, vida consagrada (871--945)}

Os fiéis de Cristo são aqueles que, pelo Baptismo, foram incorporados em Cristo e constituídos povo de Deus. Participam, cada um segundo a sua condição, do sacerdócio, do magistério e do governo de Cristo. Existe entre todos uma verdadeira igualdade quanto à dignidade e à acção comum na edificação da Igreja.

A Igreja possui uma constituição hierárquica. Cristo instituiu o ministério apostólico, confiando aos Apóstolos e aos seus sucessores o tríplice ofício de ensinar, santificar e governar. Os Bispos, em comunhão com o sucessor de Pedro, exercem este ministério pastoral.

Os leigos participam da missão da Igreja no mundo. Pelo Baptismo e pela Confirmação, são chamados a exercer o sacerdócio comum dos fiéis, oferecendo a sua vida como sacrifício espiritual e testemunhando Cristo no meio das realidades temporais.

A vida consagrada é um caminho de especial consagração a Deus através dos conselhos evangélicos. Os que professam os votos de castidade, pobreza e obediência testemunham de modo particular o Reino de Deus que está para vir.

\section{A comunhão dos santos (946--962)}

Depois de confessar a Santa Igreja Católica, o Credo acrescenta ``a comunhão dos santos''. A comunhão dos santos é a Igreja. Designa a comunhão nas coisas santas (\emph{sancta}) e a comunhão entre as pessoas santas (\emph{sancti}).

Todos os membros da Igreja estão unidos em Cristo. Os bens espirituais são comunicados entre todos. Os bens de Cristo são comunicados a todos os membros do seu Corpo. Esta comunhão existe entre os que ainda peregrinam na terra, os que se purificam no Purgatório e os que já gozam da glória do céu.

A intercessão dos santos, a oração pelos defuntos e a prática da indulgência manifestam esta comunhão viva na Igreja.

\section{Maria --- mãe de Cristo, mãe da Igreja (963--975)}

Maria, Mãe de Cristo, é também Mãe da Igreja. Ela precedeu-nos a todos no caminho da santidade. A dimensão mariana da Igreja precede a dimensão petrina. Maria é o modelo da Igreja na fé, na caridade e na união perfeita com Cristo.

Como Mãe do Redentor, tornou-se mãe de todos os membros do Corpo de Cristo. A sua maternidade estende-se a toda a Igreja. Maria é a imagem e o começo da Igreja que há-de chegar à plenitude na glória do céu. A Igreja contempla em Maria a imagem puríssima daquilo que ela própria deseja e espera ser.

A devoção à Virgem Maria, especialmente no santo Rosário e nas festas litúrgicas, é expressão da fé da Igreja na comunhão dos santos e na maternidade espiritual de Maria.